Our recent efforts to improve the reliability of spoofing countermeasures for automatic speaker verification (ASV) has confirmed that: (i)~artefacts indicative of spoofing attacks performed with converted voice or synthetic speech reside at the sub-band level~\cite{odyssey2020CQCC}; (ii)~detection performance can be improved through the use of high-spectral-resolution front-ends, particularly those which emphasize information within the same sub-bands~\cite{IS2020LFCC}.  When used even with a trivial Gaussian mixture model (GMM) back-end, high-spectral-resolution front-ends lead to extremely competitive performance, e.g.\ third position (out of 48 submissions) in the most recent ASVspoof 2019 challenge. 

Unfortunately, though, some spoofing attacks continue to evade detection, particularly the infamous A17 (.... some brief description ....) voice conversion attack within the ASVspoof 2019 database.  While we have seen that even A17 attacks can be detected when classifiers are learned with representative A17 data (i.e.\ through learning with evaluation data), when the same classifier is learned using only the official training partition, performance is poor.  This observation implies that, while high-spectral-resolution front-ends do capture A17 artefacts, classifiers cannot be learned from the available training data (i.e.\ without cheating).  The same observation shows that, once training data for a newly discovered spoofing attack is available, then any vulnerability to that attack can be fixed easily.  

There is nonetheless an interest to preempt the unexpected and to protect systems against unseen attacks.  This idea embodies the objectives of the ASVspoof initiative since its inception, namely the pursuit of generaliseable spoofing countermeasures that are capable of detecting attacks not seen in training data. The latter are characteristic of in-the-wild settings in which the nature of spoofing attacks cannot be predicted and will evolve continuously.  The question addressed in this paper, then, is whether A17 attacks can be detected even in the absence of representative training data.

Of course, this is a difficult question to answer reliably in a post-evaluation setting where the characteristics of attack A17 are now known and cannot be forgotten.  Nonetheless, there is potential for such research.  One-class classifiers trained only on bona fide data without the use of any spoofed data are an obvious candidate and have shown some previous success in anomaly, novelty and spoofing detection tasks~\cite{fatemifar2019spoofing,fatemifar2019combining,ohki2019efficient,alegreOneClass}.  Anecdotal evidence, however, shows that many, including the authors of the current article, have failed to apply them successfully to ASVspoof 2019 data.  Even though we know it not to be the case, the approach adopted here is based simply upon the assumption that A17 attacks \emph{cannot} be detected using conventional short-term spectral estimates, not even with high-spectral-resolution, sub-band features.  In this case, we must seek cues elsewhere.  The research hypothesis under investigation in this work is that, if a focus on cues localised at the spectral level are insufficient, then perhaps artefacts localised at the temporal level may offer potential.

Conventional short-term spectral estimates and GMM classifiers average scores across an utterance in its entirety and are ill-suited to the detection of temporally-localised cues.  Alternative classifier architectures with some form of temporal attention mechanism are better suited.  This paper reports our exploration of RawNet~\cite{jung2020improved} deep neural network classifiers which operate directly upon the raw speech waveform, circumventing almost entirely any use of hand-crafted features.  Identical to the SincNet~\cite{ravanelli2018speaker,ravanelli2018learning} predecessor, the first layer of the RawNet architecture is based upon a bank of sinc filters.  

Based upon the assumption that representative training data is available in abundance, in the original implementation the bandwidth and spectral position of each sinc filter is learned automatically.  Stacked convolutional and maxpooling layers... batch norm... feature map scaling... form an attention mechanism... that is capable of learning discriminative information at the spectral/temporal level.  Our assumption here is that, while a filterbank whose parameters are learned automatically may not generalise well to unseen attacks (there are only 6 attacks in the ASVspoof 2019 training data), the same architecture brings potential to learn temporally-localised artefacts.  Last, we suppose that the fusion of classifiers tuned to detect artefacts localised at both spectral and temporal levels will lead to improved performance and, perhaps optimistically, also help to detect A17 attacks.
